%Texlive-full Version 3.141592-1.40.3 (Web2C 7.5.6)
%TexStudio Version 2.7.0

\documentclass[a4paper,10pt]{article}
\usepackage[utf8]{inputenc}
\usepackage{graphicx}

\usepackage{lmodern}
\usepackage[a4paper]{geometry}

\usepackage{hyperref}

\usepackage{amsmath}
\usepackage{amssymb}
\usepackage{amsthm}

\usepackage{pstricks}
\usepackage{pst-node}
\usepackage{color}
\usepackage[ddmmyyyy]{datetime} 



\newcommand{\Bt}{B_t}
\newcommand{\RtT}{R(t,T)}
\newcommand{\PtT}{P(t,T)}
\newcommand{\PtS}{P(t,S)}
\newcommand{\LtT}{L(t,T)}
\newcommand{\YtT}{Y(t,T)}
\newcommand{\ftT}{f(t,T)}
\newcommand{\tautT}{\tau(t,T)}
\newcommand{\inttT}{\int_t^T}
\newcommand{\intt}{\int_0^t}
\newcommand{\FuncExp}{\text{exp}}
\newcommand{\FuncLn}{\text{ln}}
\newcommand{\todo}[1]{\begin{center}\color{red} TODO : #1\end{center}}

\begin{document}
\begin{titlepage}
%%%%%%%%%%%%%%%%%%%%%%%%%%  LOGO  %%%%%%%%%%%%%%%%%%%%%%%%%%%%%%%%%%%%%%%%%
\begin{center}
\begin{pspicture}(-7,0)(7,2)
\rput(-8,1){\href{http://www.univ-paris-diderot.fr/english/sc/site.php?bc=accueil&np=accueil&g=m}{\includegraphics[scale=1.2]{logo_p7}}}
\rput(4,2){\href{https://masterfinance.math.univ-paris-diderot.fr/}
	    {	\begin{tabular}{l}
		\resizebox{6cm}{2cm}{M2MO} \\
		\resizebox{6cm}{0.5cm}{DEA Laure Elie}  
		\end{tabular}
	    }
}

\rput(3,-0.5){
\begin{tabular}{l}
Tutors                                                                \\ 
\href{mailto:revierdenord@gmail.com}{\textbf{Yuan LI} }                \\
\href{mailto:scotti@math.univ-paris-diderot.fr}{\textbf{Simone SCOTTI}}   
\end{tabular}	    
}
%\psline(-7,0)(7,2)
\end{pspicture}
\end{center}



%%%%%%%%%%%%%%%%%% DOCUMENT TITLE %%%%%%%%%%%%%%%%%%%%%%%%%
\vspace{2cm}

\begin{center}
%\resizebox{14cm}{0.7cm}{High dimensional PDE solver for multi-asset pricing option}
\resizebox{14cm}{0.8cm}{Variant of Libor Market Model}
\end{center}
%%%%%%%%%%%%%%%%%% DOCUMENT TITLE %%%%%%%%%%%%%%%%%%%%%%%%%
\begin{center}
\huge{ Version : \today }
\end{center}


\begin{center}
\href{mailto:chithanhnguyen.math@gmail.com}{\LARGE{\textbf{Chi Thanh NGUYEN}}}
\end{center}



\vfill
\begin{abstract}
In the internship at Lunalogic, we continu the work of ~\cite{THAI2013} and use the book  ~\cite{BRIGO2006}
\end{abstract}
\end{titlepage}


\tableofcontents
\newpage

% % % % % % % % % % % % % % % % % % % % % % % % % % % % % % % % % % % % % % %
\newpage
\section{About Lunalogic}
\begin{tabular}{llll}
Position          &            Name             & Contact                         &Telephone      \\
Exec Director     & Fathia Rahmani               & fathia.rahmani@lunalogic.com   &               \\
HR   Director     & Isabelle Marouby             & isabelle.marouby@lunalogic.com &06.51.44.06.78 \\
Quant Consultant  & Yuan Li                      & revierdenord@gmail.com         &06.52.17.69.75 \\    
Tech. support     & Fastah Rahmani               & fastah.rahmani@gmail.com       &
\end{tabular}


% % % % % % % % % % % % % % % % % % % % % % % % % % % % % % % % % % % % % % %
\newpage
\section{Appendices}
\appendix
Some apprendices are to be added by lecture of the book ~\cite{BRIGO2006}
\subsection{Generalities about interest rate}
\begin{itemize}
 \item Zero-coupon curves, interest rate curve. When talking about curve, its is function of maturity $T$. When talking about the present time $t$, it is about process, zero-coupon process, interest rate process.
 \item Interest rate are quoted, zero-coupon are deduced from interest rate in a consistent way.
\end{itemize}

\paragraph{Money-market account (riskless asset)}
\[
d\Bt=r_t \Bt dt \hspace{1cm} B(0)=1 \hspace{1cm} B_t = \intt r_s ds
\hspace{1cm}
\left\{
\begin{array}{r l}
\Bt &\text{: Money market} \\
r_t &\text{: Instantaneous rate} 
\end{array}
\right.
\]
\paragraph{Stochastic discount factor} 
\[
D(t,T) = \FuncExp\left( -\inttT r_s ds \right)
\]
Note that $D(t,T)$ is the present value of future amount currency, while $\PtT$ zero-coupon is a contract value. $D(t,T)=\PtT$ when $r_t$ is determinist.
\paragraph{Spot Interest Rate}~\footnote{The time to maturity $\tautT$ denote the day-count convention between $t$ and $T$. The simplest case $\tautT=T-t$}\\
Continuously-compounded 
\[
R(t,T) = -\frac{\FuncLn\PtT }{\tautT}
\hspace{1.5cm}
e^{\RtT\tautT}\PtT = 1
\hspace{1.5cm}
\PtT = e^{-\RtT\tautT}
\]
Simply-compounded ~\footnote{This quantity is related to LIBOR, when time to maturity is lesser than one year} 
\[
\LtT = \frac{1-\PtT}{\tautT\PtT}
\hspace{0.7 cm}
\PtT(1+\LtT\tautT)=1
\hspace{0.7cm}
\PtT = \frac{1}{1+\LtT\tautT}
\]
Annually-compounded
\[
\YtT = \frac{1}{\PtT^{1/\tautT}} -1
\hspace{0.8cm}
\PtT(1+\YtT)^{\tautT} = 1
\hspace{0.8cm}
\PtT = \frac{1}{(1+\YtT)^{\tautT}}
\]

\paragraph{Zero-coupon curve / Yield curve / term structure of interest rate}
\[
\text{graphic of}
\hspace{2cm}
T \longmapsto
\left\{
\begin{array}{lll}
\LtT  &  t < T \leq t+1 & \text{Simply compounded} \\
\YtT  &  t+1 < T        & \text{Annually compounded} 
\end{array}
\right.
\]
Yield curve in a general way is the graph of $T \longmapsto \RtT$.
\todo{read about behavior of this curve (inverted at behavior at the begining)}
\paragraph{Zero-bond curve / term structure of discount factors}
\[
\text{Graphic of }
\hspace{2cm}
T \longmapsto \PtT
\hspace{1cm}
,t<T
\]
Note that zero-bond curve is less informative than the zero-coupon curve.
\paragraph{Forward interest rate} is the fair value of strike $K$ that make the contract \textbf{FRA} worth zero at time $t<T<S$  
\[
\text{Simply-compounded}
\hspace{2cm}
F(t,T,S) = \frac{1}{\tautT}\left( \frac{\PtT}{\PtS} - 1 \right)
\hspace{1cm}
,t<T<S
\]
Instantanuous forward interest rate ~\footnote{This fundamental quantity is related HJM model}
\[
\ftT = \lim_{S\rightarrow T^+} F(t,T,S) = -\frac{\partial \FuncLn\PtT}{\partial T}
\hspace{2cm}
\PtT = \FuncExp\left( -\inttT f(t,u)du \right)
\]

\paragraph{Swap} is a contract helping to convert fixed rate flow into floating rate flow and vice versa. 
\[
\left\{
\begin{array}{l}
T_{\alpha},T_{\alpha+1},...T_{\beta} \\
\tau_i = \tau(T_{i-1},T_i)           \\
L_i    = L(T_{i-1},T_i)
\end{array}
\right.
\hspace{2cm}
\begin{array}{l}
\text{Nominal $N$, paements at dates $T_i$,  } i=\alpha+1,...,\beta \\
\text{ - For fixed legs    } N\tau_iK  \\
\text{ - For floating legs } N\tau_iL_i  
\end{array}
\]
\paragraph{Caps / Floors}  
\[
\left\{
\begin{array}{l}
T_{\alpha},T_{\alpha+1},...T_{\beta} \\
\tau_i = \tau(T_{i-1},T_i)           \\
L_i    = L(T_{i-1},T_i)
\end{array}
\right.
\hspace{2cm}
\begin{array}{l}
\text{Nominal $N$, payments at dates $T_i$,  } i=\alpha+1,...,\beta \\
\text{ - For Caps    } N\tau_i(L_i-K)^+                              \\
\text{ - For Floors  } N\tau_i(K-L_i)^+
\end{array}
\]
\paragraph{Swaption}
\todo{Black Formula for pricing Caps, Floors, Swaptions}  









relation between monney market account, zerocoupon, spot rate, short rate, forward rate , swaps, caps, floor, swaptions... 
\subsection{Forward price, change of numeraire, Forward measure}
Definitions of forward price, forward measure, change of numeraire. Pricing Caps, Floors, Swap?, Swaption?

\subsection{HJM Framework}
\subsection{Correlated Brownian motion, volatility}
\subsection{LIBOR Market Model}
\subsection{Software and systems}
Installed on PC : TortoiseSVN, tex-live, texstudio, nodepadd++, Visual Studio Express 2012, cmake, boost\_1\_55\_0 (self full built), QuantLib-1.4, BOUML, 7-zip, GLE, GSview. 

% % % % % % % % % % % % % % % % % % % % % % % % % % % % % % % % % % % % % % %
\newpage
\bibliographystyle{siam}
\bibliography{intern_lunalogic.bib}
\end{document}
